% Generated draft: requires mathematical and source review.
\documentclass[11pt]{article}
\usepackage[margin=25mm]{geometry}
\usepackage{fontspec}
\setmainfont{Latin Modern Roman}
\usepackage{amsmath,amssymb,mathtools,unicode-math}
\setmathfont{Latin Modern Math}
\usepackage{enumitem,needspace,xurl,hyperref}
\setlength{\parindent}{0pt}
\setlength{\parskip}{5pt}
\setlength{\emergencystretch}{3em}
\newcommand{\N}{\mathbb{N}}
\newcommand{\Z}{\mathbb{Z}}
\newcommand{\Q}{\mathbb{Q}}
\newcommand{\R}{\mathbb{R}}
\newcommand{\C}{\mathbb{C}}
\newcommand{\F}{\mathbb{F}}
\newcommand{\A}{\mathbb{A}}
\newcommand{\PP}{\mathbb{P}}
\newcommand{\E}{\mathbb{E}}
\newcommand{\eps}{\varepsilon}
\newcommand{\dd}{\,\mathrm{d}}
\begin{document}
\section*{UnsolvedMath ID: 30003466}
\textbf{Source problem code:} catalogue entry
\subsection{Definitions and mathematical statement}
Under the standard convex-body interpretation of the supplied statement, a plane convex body is a compact convex set \(D\subset\R^2\) with nonempty interior. A positive homothetic copy of \(D\) is
\[
t+\lambda D=\{t+\lambda x:x\in D\},\qquad t\in\R^2,\quad\lambda>0.
\]
For a finite set \(P\subset\R^2\), a three-coloring is a map \(c:P\to\{1,2,3\}\); it need not use all three colors. A subset of \(P\) is nonmonochromatic if it contains two points with distinct colors.

\textbf{Question under these conventions.} Is it true that for every plane convex body \(D\) there exists an integer \(m=m(D)\ge2\) such that, for every finite set \(P\subset\R^2\), there exists a coloring \(c:P\to\{1,2,3\}\) satisfying
\[
\forall t\in\R^2\ \forall\lambda>0,\qquad
\bigl|P\cap(t+\lambda D)\bigr|\ge m
\ \Longrightarrow\
\exists p,q\in P\cap(t+\lambda D)\quad c(p)\ne c(q)?
\]
Here \(m\) depends only on \(D\), and the same coloring must work simultaneously for every indicated copy. Because \(D\) is closed under the stated convex-body convention, points on the boundary of a copy count toward its intersection with \(P\).

\textbf{Unresolved source qualifications.} The supplied cleaned statement says ``plane convex body,'' whereas the original statement says ``plane convex set.'' The displayed question follows the cleaned wording with the standard definition above; the original wording does not expressly impose compactness, nonempty interior, or closedness. The supplied record also does not specify whether homothetic copies allow negative scale factors or how boundary points should be treated independently of the convex-body convention. These scope and convention choices require source review.
\subsection{Short English statement}
For each fixed plane convex body, ask whether a threshold exists such that every finite planar point set can be colored with three colors so that every positive homothetic copy containing at least that many points contains two different colors. The threshold must be independent of the point set. The original wording uses the broader term ``convex set,'' so the intended domain and geometric conventions require review.
\subsection{Sources}
\begin{itemize}
\item Supplied local UnsolvedMath registry record, permanent ID \(30003466\), title ``Proper Three-Colorings for Homothets of Convex Sets,'' source problem code \texttt{catalogue entry}, in the Oberwolfach Reports Open Problems collection. No individual proposer is named; the supplied proposed year is 2017.
\item Supplied report citation: \emph{Discrete Geometry (2018)}, Oberwolfach Reports, DOI 10.4171/owr/2017/19, pp.\ 1135--1205. \url{https://doi.org/10.4171/owr/2017/19}. The supplied background describes Oberwolfach Reports as documenting mathematical workshops and their open-problem discussions. The citation's 2018 label and the record's 2017 year indicators are retained without reconciliation.
\item Supplied literature source: Keszegh and P\'alv\"olgyi on proper coloring, \url{https://arxiv.org/abs/1612.02158}.
\item Supplied literature source labeled ``original OWR report'': \url{https://doi.org/10.4171/owr/2017/19}.
\item Supplied dated literature curation, checked \texttt{2026-08-21}: the record labels the problem open and reports that the disk case remains open, while the cited primary paper proves the analogous statement for convex polygons but not for disks. This assessment is supplied dated curation derived from the OWR curation, not independently verified here; its time-sensitive status and cited claims require independent verification.
\item The registry labels the statement \texttt{corrected\_verified} and describes repairs to notation, typography, and sentence structure while preserving mathematical content. These are supplied editorial assertions, not independent mathematical verification of this draft.
\end{itemize}
\end{document}
